\section{Limitations}
\label{appendix:limitations}

While \benchname\space provides a comprehensive evaluation framework for RAG systems, it has a few limitations that should be acknowledged and addressed in future research.

First, the diagnostic metrics for the retriever component are less insightful compared to those for the generator. The retrieval metrics primarily focus on the recall of ground truth claims and precision of retrieved context, but they may not fully capture the nuances and complexities of the retrieval process. Developing more sophisticated metrics that consider factors such as the information density, diversity and coherence of the retrieved context could provide deeper insights into the retriever's performance.


Second, the metrics proposed in \benchname\space do not differentiate between \texttt{Neutral} and \texttt{Contradiction} checking results from RefChecker when evaluating the generated responses. These two types of results may have different impacts on the final response quality, and treating them equally could lead to an incomplete assessment. Future work should explore ways to incorporate the distinction between neutral and contradiction results into the evaluation metrics, potentially assigning different weights or penalties based on their severity.

Finally, the evaluation benchmark used in this study is curated based on existing text-only datasets and is limited to English queries and corpus. While this allows for a focused evaluation of RAG systems, it may not fully represent the diverse range of tasks and languages that RAG systems can be applied to. Expanding the benchmark to include datasets from different modalities (e.g., images, audio) and languages would provide a more comprehensive assessment of RAG systems' capabilities and generalization. Additionally, creating benchmark datasets specifically designed for evaluating RAG systems, rather than repurposing existing ones, could help to better capture the unique challenges and requirements of this task.

By refining the diagnostic metrics, incorporating the impact of different checking results, and expanding the evaluation benchmark, researchers can gain an even more comprehensive understanding of RAG systems' performance and identify targeted areas for improvement.